% F7 -- Life cycle of a branch: the state machine of Definition "Branch state".
% Requires hpb.sty v1.0 (palette and glyph styles). Plain TikZ, units in cm.
\begin{figure}[t]
\centering
\resizebox{\linewidth}{!}{%
\begin{tikzpicture}[
    x=1cm, y=1cm,
    font=\sffamily\fontsize{7}{8.4}\selectfont,
    lc/state/.style   = {rectangle, rounded corners=5pt, draw=hpbInk,
                         line width=1pt, minimum width=3.3cm,
                         minimum height=1.45cm, inner sep=0pt},
    lc/final/.style   = {double=white, double distance=1.3pt, line width=0.8pt},
    lc/name/.style    = {font=\sffamily\bfseries\fontsize{8}{9.6}\selectfont,
                         text=hpbInk, anchor=west, inner sep=0pt},
    lc/inv/.style     = {font=\sffamily\fontsize{6.5}{7.8}\selectfont,
                         text=hpbMuted, anchor=west, inner sep=0pt},
    lc/arrow/.style   = {draw=hpbInk, line width=1pt, rounded corners=6pt,
                         -{Stealth[length=5pt,width=4.5pt]}},
    lc/audit/.style   = {lc/arrow, dash pattern=on 4pt off 2.5pt},
    lc/tr/.style      = {font=\sffamily\fontsize{7}{8.4}\selectfont,
                         text=hpbInk, fill=white, inner sep=1.5pt},
    lc/trb/.style     = {lc/tr, font=\sffamily\bfseries\fontsize{7}{8.4}\selectfont},
  ]
  % ---- the open set W(t): states in which a branch counts as work in progress
  % (dotted frame: deliberately not the style of an HPB gate band)
  \draw[draw=hpbDiscInk, line width=0.8pt, rounded corners=4pt,
        dash pattern=on 0.1pt off 2pt, line cap=round]
    (1.45,-1.02) rectangle (10.15,1.02);
  \node[lc/inv, anchor=south west, font=\sffamily\itshape\fontsize{6.5}{7.8}\selectfont]
    at (1.5,1.06) {open set $W(t)$};

  % ---- states: neutral fills (not the baseline color of the notation);
  %      the glyph in each box is the HPB terminal that represents the state
  \node[lc/state, fill=white]               (open) at (3.35,0)  {};
  \node[lc/state, fill=white]               (targ) at (8.25,0)  {};
  \node[lc/state, lc/final, fill=white]     (merg) at (13.65,0) {};
  \node[lc/state, lc/final, fill=white]     (disc) at (8.25,-3) {};

  % glyphs (shown for an improvement)
  \node[hpb/head ring, draw=hpbImp] at (2.2,0) {};
  \node[hpb/event, fill=white]      at (2.2,0) {};
  \draw[hpb/path, hpb/path target] (7.3,0) -- (7.83,0);
  \node[hpb/head ring, draw=hpbImp] at (7.1,0) {};
  \node[hpb/event, fill=white]      at (7.1,0) {};
  \node[hpb/event, fill=hpbImp]     at (12.5,0) {};
  \draw[draw=hpbInk, line width=1.3pt, line cap=round, line join=round]
    (12.5,0) ++(-2.6pt,0.1pt) -- ++(1.8pt,-2pt) -- ++(3.5pt,4.2pt);
  \node[hpb/event, fill=hpbDisc, draw=hpbMuted] at (7.1,-3) {};
  \draw[draw=hpbInk, line width=1.3pt, line cap=round]
    (7.1,-3) ++(-2.2pt,-2.2pt) -- ++(4.4pt,4.4pt) ++(-4.4pt,0) -- ++(4.4pt,-4.4pt);

  % names and invariants
  \node[lc/name] at (3.0,0.22)   {Open};
  \node[lc/inv]  at (3.0,-0.2)   {$g(b)$ unset};
  \node[lc/name] at (7.95,0.22)  {Targeted};
  \node[lc/inv]  at (7.95,-0.2)  {$g(b)$ planned;};
  \node[lc/inv]  at (7.95,-0.47) {verified at $t_r(b)$};
  \node[lc/name] at (13.4,0.22)  {Merged};
  \node[lc/inv]  at (13.4,-0.2)  {$m(b)=g(b)$};
  \node[lc/name] at (7.95,0.22-3) {Discarded};
  \node[lc/inv]  at (7.95,-0.2-3) {rationale $\rho(b)$};

  % ---- creation: a branch is opened by forking from a released baseline,
  %      drawn as the released-baseline glyph of the notation
  \node[hpb/baseline] (init) at (-0.35,0) {};
  \draw[lc/arrow] (init.east) -- (open.west)
    node[lc/trb, midway, above=2pt] {fork}
    node[lc/inv, midway, below=3pt, anchor=north] {from $f(b)$};

  % ---- ordinary transitions
  \draw[lc/arrow] ([yshift=3.5mm]open.east) -- ([yshift=3.5mm]targ.west)
    node[lc/trb, midway, above=1pt] {target};
  \draw[lc/arrow] ([yshift=-3.5mm]targ.west) -- ([yshift=-3.5mm]open.east)
    node[lc/trb, midway, below=1pt] {un-target};
  \draw[lc/arrow] (7.75,0.725) -- (7.75,1.3) -- (8.75,1.3) -- (8.75,0.725);
  \node[lc/trb, anchor=south] at (8.25,1.32) {re-target};
  \draw[lc/arrow] (targ.east) -- (merg.west)
    node[lc/trb, midway, above=2pt] {merge}
    node[lc/inv, midway, below=3pt, anchor=north, align=center]
      {on publication\\of $g(b)$, if verified};
  \draw[lc/arrow] ([xshift=6mm]open.south) |- (disc.west)
    node[lc/trb, pos=0.75, above=1pt] {discard};
  \draw[lc/arrow] (targ.south) -- (disc.north)
    node[lc/trb, midway, right=2pt] {discard};

  % ---- exceptional, audited transitions (dashed)
  \draw[lc/audit] (merg.north) -- (13.65,1.95) -- (3.95,1.95) -- (3.95,0.725);
  \node[lc/trb, anchor=south] at (8.8,1.97) {reopen\quad
    \textmd{\textcolor{hpbMuted}{(exceptional, audited)}}};
  \draw[lc/audit] (disc.south) -- (8.25,-4.15) -- (2.75,-4.15) -- (2.75,-0.725);
  \node[lc/trb, anchor=north] at (5.5,-4.17) {reopen\quad
    \textmd{\textcolor{hpbMuted}{(exceptional, audited)}}};
\end{tikzpicture}}
\caption{Life cycle of a branch (\cref{def:state}). A branch is
\textsf{open} until it is committed to a planned version (\textsf{targeted});
it merges only from \textsf{targeted}, when that version is published and
after it has been verified; it can be discarded from either state, always
with a reason. \textsf{Merged} and \textsf{discarded} are final; dashed
transitions are exceptional and audited. Each box shows the glyph of the
state; the dotted frame marks the states counted as work in progress,
$W(t)$.}
\label{fig:lifecycle}
\end{figure}
